\subsection{Olas, perfiles y reglas de habilitación}
La Tabla~\ref{tab:t18-olas-perfiles} vincula las cuatro olas de cada etapa con procesos, zonas y perfiles. E1 se habilitará del 1 al 14 de diciembre de 2027; E2, en junio de 2028. En E1 el varadero sólo recibe preparación y acompañamiento del procedimiento vigente: su alcance nuevo se habilita en E2. La habilitación del alcance obligatorio completo precede a sus cuatro semanas finales de observación. Cada acta de ola identificará zona, versión, perfiles habilitados, fecha, autorizador y controles CA aplicables; una ola no autoriza activar una función de otra etapa.
\begin{longtable}{L{10mm}L{42mm}L{60mm}L{37mm}}
\caption{Olas, zonas y perfiles de implantación}\label{tab:t18-olas-perfiles}\\\hline
 \EncabezadoTabla{Ola} & \EncabezadoTabla{Procesos y zonas} & \EncabezadoTabla{Perfiles} & \EncabezadoTabla{Responsabilidad}\\\hline\endfirsthead
\hline \EncabezadoTabla{Ola} & \EncabezadoTabla{Procesos y zonas} & \EncabezadoTabla{Perfiles} & \EncabezadoTabla{Responsabilidad}\\\hline\endhead
1 & Pantalanes, torre y zarpes & Marineros, torre, portería; avanzados y soporte & Implantación; validan Operaciones y Calidad\\\hline
2 & Escuela y administración & Monitores, administración y referentes técnicos; apoderados en su portal & Implantación; validan Escuela y Administración\\\hline
3 & Varadero, combustible y medición & Operadores de izaje, surtidor y contratistas; soporte & Implantación; validan responsables operacionales\\\hline
4 & Portal y coordinación transversal & Armadores, apoderados, contratistas, administradores y soporte & Implantación; validan propietarios de proceso\\\hline
\end{longtable}
\FuenteTabla{Elaboración propia: plan de implantación por procesos, etapas y criterios CA de SD3; condiciones de cierre del artículo 17.3 de las Bases Administrativas, p. 13.}
José Lara Arce coordinará implantación y referentes; Davor Serey verificará evidencia y preparación; Ignacio Silva coordinará la decisión con la Contraparte. El CLIENTE designará por proceso un referente y un suplente antes de habilitar la ola. No se supone área TI interna ni formador propio: Synaptix ejecutará formación, soporte y mentoría; los referentes de la marina validarán tareas, decisiones y prioridades de negocio.

\subsection{Capacitación, evaluación y preparación por sitio}
Cada etapa mantiene 107 cohortes de hasta doce plazas por rol y 1.284 HH de proveedor: cien de usuarios finales, dos de avanzados, una de administradores, una de responsables técnicos designados y tres de soporte N1/N2. Cada cohorte incluye preparación, dos sesiones de dos horas, ejercicio, evaluación individual y recuperación; se imputan Implantación 8, Funcional 2 y Calidad 2 HH. Las cuatro horas de sesión por participante no se confunden con las doce HH del proveedor. Los registros \texttt{1.11.7.e.p.s.b} identifican sitio, turno y cohorte; la lista nominal asignará cada persona a sus roles antes de las sesiones.
Se combinarán presencialidad en cada sitio operativo, sesiones sincrónicas, autoformación en línea con avance registrado y acompañamiento en puesto durante marcha blanca. Los materiales por perfil incluirán manual editable, guía rápida, FAQ, video tutorial en español y base de conocimiento, de propiedad del CLIENTE y vinculados a versión. Las familias \texttt{1.11.2.e.p} producen materiales; \texttt{1.11.3.e.1} fija el plan por turno/sitio y \texttt{1.11.3.e.2} el expediente de evaluación y certificación. La asistencia sin ejercicio aprobado no acredita certificación; administradores y referentes técnicos deben estar certificados antes de cerrar cada marcha blanca.
La formación E1 termina el 24 de noviembre de 2027 y la E2 el 11 de mayo de 2028, conforme a las puertas FORMACION-E1/E2. Las recuperaciones obligatorias se completarán antes de habilitar el perfil y fuera de congelamiento; Synaptix añadirá capacidad si resultara necesaria. La incorporación estacional reserva 25 cohortes en noviembre de cada ciclo, además de dos jornadas anuales de actualización y de la formación sin recargo del personal nuevo. Monitores nuevos y suplentes se incluirán expresamente, sin depender de un instructor interno.~\parencite[arts.~89--90; p.~45]{bases-admin}~\parencite[RT-22.01--06; p.~38]{bases-transversales}

\subsection{Adopción y acompañamiento decreciente}
El plan \texttt{1.11.3.e.1} incorporará diagnóstico por área/perfil, resistencias, agentes de cambio designados y suplentes, calendario/canales de comunicación y procedimientos actualizados. Los ejercicios de pantalán se realizarán con guantes y condiciones seguras de humedad; varadero ensayará el registro antes y después de la maniobra, nunca con carga suspendida. El canal voluntario de armador mantendrá consentimiento, revocación y alternativa funcional sin adhesión. La formación no sustituye la comprobación de estas condiciones de uso.
Por ola y zona se reserva un puesto de ocho horas de presencia durante los primeros catorce días. En los siguientes catorce, se propone cuatro horas presenciales más cuatro de preparación, seguimiento y evidencia: 448 HH por ola, 1.792 por etapa. La reducción sólo se autorizará tras siete días consecutivos con al menos 90\,\% de usuarios activos sobre usuarios con tarea exigible, 95\,\% de ejercicios/tareas sin asistencia y 100\,\% de actuaciones críticas correctas, sin incidente crítico/alto abierto. Son metas de oferta; denominadores excluyen personas sin turno o sin tarea en el período, con exclusión justificada. No se equiparan las plazas de formación a usuarios activos.
Calidad medirá diariamente cobertura, uso por función, abandono de mecanismos manuales no exigidos, necesidad de ayuda y resultados críticos; Implantación recogerá satisfacción semanal, con meta de 80\,\% favorable entre respuestas válidas y tasa de respuesta declarada. Ante rechazo, dificultad o meta incumplida, se mantiene o refuerza presencia, se adapta el procedimiento y se programa recuperación permitida; Synaptix financia la capacidad adicional sin detraer la cobertura de servicio. La permanencia del legado durante marcha blanca no se contabiliza como fracaso de adopción. La Contraparte autorizará reducción y avance con el expediente de ola.~\parencite[art.~89; p.~45]{bases-admin}~\parencite[RT-20.01/05; p.~35]{bases-transversales}

\subsection{Cuadro diario de control y cierre}
La Tabla~\ref{tab:t18-indicadores} vincula los indicadores diarios con su evidencia y acción de control. Cada proceso conservará su denominador de hechos reales, cantidades absolutas, mezcla y períodos sin actividad. El cuadro aprobado por proceso conservará como piso los mínimos absolutos y las fórmulas de la Tabla~\ref{tab:sd7-volumen-real}, con mezcla y calendario real; no basta cumplir un porcentaje sobre un volumen insuficiente. Ninguna carga sintética ni porcentaje sobre procesos parcialmente habilitados sustituirá el alcance completo exigible de la etapa.
\ifdefined\EnCuerpoSD
\begin{table}[H]
\centering
\begin{tabularx}{\linewidth}{L{38mm}Y L{32mm}}
\hline
 \EncabezadoTabla{Indicador} & \EncabezadoTabla{Umbral de cierre} & \EncabezadoTabla{Responsable de control}\\\hline
Cobertura de hechos reales & 100\,\% durante 28 días finales & Datos y Calidad\\\hline
Conciliación & Cero diferencias no explicadas, pérdidas o duplicados & Datos y Funcional\\\hline
Disponibilidad y tiempos & Umbrales CA/SLA de SD3--SD5 durante 28 días & Operación y Calidad\\\hline
Incidentes y formación & Cero críticos/altos abiertos; perfiles preparados y certificación exigida & Calidad e Implantación\\\hline
Acta y autoridad & Firma expresa anterior al cambio oficial & Proyecto\\\hline
\end{tabularx}
\caption{Síntesis de indicadores diarios y condiciones de cierre}
\label{tab:t18-indicadores}
\FuenteTabla{Elaboración propia: plan de implantación por procesos, etapas y criterios CA de SD3; condiciones de cierre del artículo 17.3 de las Bases Administrativas, p. 13.}
\end{table}
\else
\begin{longtable}{L{35mm}L{69mm}L{45mm}}
\caption{Indicadores de avance y cierre}\label{tab:t18-indicadores}\\\hline
 \EncabezadoTabla{Indicador} & \EncabezadoTabla{Fuente, denominador y umbral} & \EncabezadoTabla{Control y acción}\\\hline\endfirsthead
\hline \EncabezadoTabla{Indicador} & \EncabezadoTabla{Fuente, denominador y umbral} & \EncabezadoTabla{Control y acción}\\\hline\endhead
Volumen y cobertura & Diario del registro vigente frente a manifiestos; 100\,\% de hechos exigibles de todos los procesos de la etapa durante 28 días finales; cantidades aprobadas por proceso & Datos y Calidad; una omisión impide cerrar\\\hline
Conciliación & Recuentos, sumas por moneda/período, documentos y muestreo dirigido de SD5; cero diferencias no explicadas, cero pérdidas o efectos duplicados & Datos y validador funcional; detener grupo afectado y corregir/retornar\\\hline
Disponibilidad y tiempos & Telemetría y transacciones completas, intervalos y exclusiones declarados; umbrales CA de SD3 y SLA de SD4/SD5 sostenidos los mismos 28 días & Operación y Calidad; incumplimiento impide cierre\\\hline
Incidentes y formación & Registro de severidad; cero críticos/altos abiertos atribuibles; 100\,\% de perfiles obligatorios preparados y certificación de administradores/técnicos & Calidad e Implantación; recuperar evidencia antes del acta\\\hline
Acta y autoridad & Expediente con versión, evidencia y observaciones resueltas; firma expresa de Contraparte previa al cambio oficial & Proyecto; sin acta no habilitar registro oficial\\\hline
\end{longtable}
\FuenteTabla{Elaboración propia: plan de implantación por procesos, etapas y criterios CA de SD3; condiciones de cierre del artículo 17.3 de las Bases Administrativas, p. 13.}
\fi

Los criterios CA se aplican por proceso y etapa, sin sustituir umbrales por promedios de toda la solución. Las actas y protocolos se formalizarán en T-17; la definición de terminado incluye resultado/código, pruebas, documentación, seguridad y despliegue. Durante marcha blanca el registro oficial vigente conserva autoridad; la nueva solución se usa con usuarios/datos reales y escrituras supervisadas para conciliación. La habilitación oficial E1 será en mes 16; E2, desde el primer día de 21 con acta previa.~\parencite[arts.~17--18; pp.~12--13]{bases-admin}~\parencite[RT-20.04/07/08; p.~35]{bases-transversales}
